%----------------------------------------------------------------------------------------
%	PACKAGES AND OTHER DOCUMENT CONFIGURATIONS
%----------------------------------------------------------------------------------------
\DocumentMetadata{
pdfversion=2.0,  
lang=en-US,   
pdfstandard={ua-2, a-4f},
tagging = on, 
tagging-setup={math/setup={mathml-SE}} ,
}
\documentclass[12pt]{article}
\usepackage{amsmath}
\usepackage[extreme]{savetrees}
\usepackage[utf8]{inputenc}
\usepackage{newpxtext}
\usepackage{hyperref}
\usepackage{multicol}
\usepackage{physics}
\usepackage{siunitx}

%%%%%%%%% === Document Configuration === %%%%%%%%%%%%%%

\hypersetup{
pdftitle={MATH 340 -- Homework 7 Spring 2026},
pdfauthor={Ignacio Rojas},
pdfkeywords={seventh homework},%
}

\newcommand{\twobyone}[2]{\begin{pmatrix} %% 2 x 1 matrix
   #1 \\ #2 \end{pmatrix}}
\newcommand{\twobytwo}[4]{\begin{pmatrix} %% 2 x 2 matrix
   #1 & #2 \\ #3 & #4 \end{pmatrix}}
\newcommand{\threebyone}[3]{\begin{pmatrix} %% 3 x 1 matrix
   #1 \\ #2 \\ #3 \end{pmatrix}}
\newcommand{\threebythree}[9]{\begin{pmatrix} %% 3 x 3 matrix
   #1 & #2 & #3 \\ #4 & #5 & #6 \\ #7 & #8 & #9 \end{pmatrix}}

%----------------------------------------------------------------------------------------
%	ARTICLE CONTENTS
%----------------------------------------------------------------------------------------
\begin{document}

\begin{center}
    {\Large MATH 340 -- Homework 7,\quad Spring 2026}
\end{center}

Recall that you must hand in a subset of the problems for which deleting any problem makes the total score less than 10. The maximum possible score on this homework is 10 points. See the syllabus for scoring details. NOTE: All problems from Boyce, Di Prima and Meade are from the \(11^{\text{th}}\)  edition, which is available online: \href{https://www.math.colostate.edu/~liu/MATH340_S26Crd/Textbook_BoyceDiPrimaMeade_11ed.pdf}{https://www.math.colostate.edu/\(\sim\)liu/MATH340\_S26Crd/Textbook\_BoyceDiPrimaMeade\_11ed.pdf}.

\begin{enumerate}
  
  \item (2+) [4 points] For a real constant \(a\), consider the systems of equations 
  \[
  \left\lbrace
  \begin{aligned}
   &ax+y+2z=0,\\
   &x+y+z=0,\\
   &-x+y+(1-a)z=0.
  \end{aligned}
  \right.
  \]
  \begin{enumerate}
   \item Build the associated matrix \(A\) to the system of equations.
   \item Find the values of \(a\) such that \(\det A=0\).
   \item Row reduce the matrices you obtain for such values of \(a\).
   \item Write down the solution to the system of equations in such cases.
  \end{enumerate}

  \item (2) [3 points] For a real constant \(k\), consider the augmented matrix
  \[
  \left(\begin{array}{ccc|c}
      1&k & -1 & 1 \\
      0 & -2 & k+1 & 2k  \\
      0 & 0 & k^2+5k+6 & k+3  
    \end{array}\right)
    \]
    \begin{enumerate}
      \item Write down the associated system of equations associated to this augmented matrix.
      \item For what value of \(k\) does the previous system have no solution?
      \item Write down what ``the equation which causes the problem'' would look like for the previously found value of \(k\)
    \end{enumerate}

    \item (2) [3 points] \textbf{The 2\(\times\) 2 eigenvector trick}: Consider a \(2\times 2\) matrix \(A\) with an eigenvalue \(\lambda\). Verify that if 
    \[A-\lambda I=\twobytwo{z}{w}{\ast}{\ast}\]
   then the vector \((-w,z)\), is an eigenvector of \(A\) associated to the eigenvalue \(\lambda\). [If the row \(\begin{pmatrix}z&w\end{pmatrix}\) is zero, just use the other row.]\par
   The \textbf{only case} where this doesn't work is if originally your matrix was 
   \[
   A=\twobytwo{a}{0}{0}{a},\quad\text{for a constant }a.
   \]
   No need to answer, but, think: what happens on that case?

   \item (3-) [8 points] \textbf{The (other) 2\(\times\) 2 eigenvector trick}: Again, suppose \(A\) is a \(2\times 2\) matrix with eigenvalues \(\lambda_1,\lambda_2\) (possibly equal). Verify that if 
    \[A-\lambda_1 I=\twobytwo{r}{\ast}{s}{\ast}\]
    then the vector \((r,s)\) is an eigenvector of \(A\) associated to the \emph{other} eigenvalue \(\lambda_2\). [Again, if the column in question is zero, use the other one.]\par
    As before, this \textbf{doesn't work} if your matrix was originally
    \[
   A=\twobytwo{a}{0}{0}{a},\quad\text{for a constant }a.
   \]
   What happens for such diagonal matrices?
   

\item (1+) [2 points] Boyce, di Prima, and Meade chapter 7 section 5 problems 2 and 4.
  
\item (1-) [1 points] Boyce, di Prima, and Meade chapter 7 section 5 problem 11.

\item (1+) [2 points] Boyce, di Prima, and Meade chapter 7 section 6 problems 3 and 4.

\item (1-) [1 points] Boyce, di Prima, and Meade chapter 7 section 6 problem 7.

\item (1+) [2 points] Boyce, di Prima, and Meade chapter 7 section 8 problems 2 and 3.

\item (1-) [1 point] Boyce, di Prima, and Meade chapter 7 section 8 problem 6.

\item (1+) [2 points] Boyce, di Prima, and Meade chapter 7 section 5 problem 18.

\item (1+) [2 points] Boyce, di Prima, and Meade chapter 7 section 6 problem 9.

\item (3-) [8 points] Boyce, di Prima, and Meade chapter 7 section 6 problem 25. 

\item (\textbf{E}) (1+) [2 points] Consider the system \(\dot{x}=Ax\) with \(A=\twobytwo{0}{1}{1}{0}\). Find the specific solution to this system satisfying \(x_p(0)=(4,2)\).

\item (\textbf{E}) (2) [3 points] Consider the system \(\dot{x}=Ax\) with \(A=\threebythree{0}{1}{0}{-2}{2}{0}{0}{0}{-3}\). Find the \emph{real-valued} general solution to this system.

\item (\textbf{E}) (1+) [2 points] Consider the system \(\dot{x}=Ax\) with \(A=\twobytwo{2}{4}{-1}{-2}\). If \(x_p\) is specific solution satisfying \(x_p(0)=(3,2)\), find \(x_p(1)\).

\item (\textbf{E}) (2) [3 points] Consider the system \(\dot{x}=Ax\) with \(A=\threebythree{-2}{0}{1}{0}{1}{0}{0}{0}{-2}\). Find the general solution to this system.

\item (\textbf{E}) (1-) [1 point] The nonlinear system
\[
\left\lbrace
\begin{aligned}
   &\dot{x}(t)=\sin(y),\\
   &\dot{y}(t)=x,
\end{aligned}
\right.
\]
has an equilibrium point at \((0,\pi)\). What is that equilibrium point's behavior?

\item (\textbf{E}) (2) [3 points] Consider the nonlinear system
\[
\left\lbrace
\begin{aligned}
   &\dot{x}(t)=-2y+4,\\
   &\dot{y}(t)=-3x^2+12,
\end{aligned}
\right.
\]
Find its equilibria and classify them by stability type.


\end{enumerate}

\end{document}